\chapter{Background}\label{ch:background}

This chapter provides the technical foundation for the approach presented in
\cref{ch:methods}.
It introduces \sysml and its Interaction Model constructs. It explains the property
graph data model. It describes the static and dynamic analysis techniques used.
Finally, it summarises the dependence-based slicing methodology.
This chapter contains background knowledge rather than contributions.

\section{SysML v2 and Interaction Models}\label{sec:sysml}

The \acl{OMG} \acl{SysML} version~2 (\sysml) is a general-purpose modelling language
for systems engineering. It is built on top of the Kernel Modeling Language (\ac{KerML})
\cite{omg:sysmlv2}.
Unlike its predecessor (SysML v1, which was a profile of UML), \sysml defines
its own textual and graphical notation. This enables model interchange via a
standardised textual syntax.

\subsection{Sequence Diagrams in SysML v2}

Interaction Models in \sysml are expressed as \emph{occurrence specifications}
that describe ordered sequences of message occurrences between \emph{parts}
(\ie instances of blocks in the system structure).
The constructs used by the generator in this thesis (\cref{sec:sysml-impl}) are:
\begin{description}
  \item[\texttt{part def}~$:>$~\texttt{InteractionScenario}] A named
    container for a scenario-specific interaction model, analogous to a
    \ac{UML} interaction.
  \item[\texttt{Lifeline}] A participant in the interaction, \ie an
    instance of a structural block along which message occurrences are
    ordered.
  \item[\texttt{SynchronousCall} / \texttt{ReturnMessage} /
    \texttt{AsynchronousMessage}] Message stereotypes (specialising a
    common \texttt{Message} definition) that model a call, its return, and
    an asynchronous send between two lifelines.
  \item[\texttt{from} / \texttt{to}] References on a message that bind it
    to its sending and receiving lifelines, giving each message occurrence
    its position in the ordered sequence.
\end{description}

An example in \sysml textual notation, corresponding to the S1
scenario (Aggregator notification chain) used in this thesis and following
the same stereotype-based form the generator emits (\cref{sec:sysml-impl}):

\begin{minted}[fontsize=\footnotesize]{text}
package AggregatorNotification {
    part def SynchronousCall :> Message { ... }

    part def S1_AggregatorChain :> InteractionScenario {
        part runner : Lifeline;
        part agg    : Lifeline;
        part log    : Lifeline;

        part m1 : SynchronousCall {
            ref from : Lifeline = runner;
            ref to   : Lifeline = agg;
            attribute :>> label = "1. getAll";
        }
        part m2 : SynchronousCall {
            ref from : Lifeline = runner;
            ref to   : Lifeline = log;
            attribute :>> label = "2. stream";
        }
        part m3 : SynchronousCall {
            ref from : Lifeline = runner;
            ref to   : Lifeline = log;
            attribute :>> label = "3. flush";
        }
        part m4 : SynchronousCall {
            ref from : Lifeline = runner;
            ref to   : Lifeline = log;
            attribute :>> label = "4. instance";
        }
    }
}
\end{minted}

\section{Property Graphs}\label{sec:property-graphs}

A \acl{PG} is a directed, attributed multigraph in which both nodes and
edges carry arbitrary key-value properties \cite{angles:pg-survey}.
Formally, a \ac{PG} is defined as $G = (V, E, \lambda, \mu)$ where:
\begin{itemize}
  \item $V$ is the set of nodes in the graph.
  \item $\lambda(v) \in \mathcal{L}$ assigns a label from the label set
    $\mathcal{L}$ to each node $v \in V$.
  \item $E \subseteq V \times \mathcal{R} \times V$ is the set of directed edges.
    Each edge carries a relation type $r$ drawn from the relation-type set
    $\mathcal{R}$.
  \item $\mu : (V \cup E) \times \mathcal{K} \rightharpoonup \mathcal{V}$
    is a partial function that assigns property values from value set
    $\mathcal{V}$ to nodes or edges, indexed by property keys from key set
    $\mathcal{K}$.
\end{itemize}

\neo is the most widely deployed \ac{GDB} that implements the \ac{PG} model.
It is queried using the Cypher declarative query language \cite{neo4j:cypher}.
The choice of a \ac{PG} over a relational or document store is motivated by the
need to traverse heterogeneous analysis artefacts (functions, files, modules,
and trace events) efficiently along typed relationships.

\section{Static Call-Graph Analysis}\label{sec:static-analysis}

A \emph{call graph} is a directed graph in which nodes represent functions
(or methods) and edges represent potential call relationships
\cite{grove:callgraph}.
An edge $(f, g)$ exists if and only if function $f$ may call function $g$
under some execution. Here, $f$ denotes the caller function and $g$ denotes
the callee function.

\subsection{Clang and the Renaissance Semantic Extractor}

Clang is the \Cpp\ frontend of the LLVM compiler infrastructure. It exposes
a rich \ac{AST} traversal \ac{API} that enables static analysis tools to
extract precise call-graph information \cite{llvm:clang}.
Rather than driving Clang directly, this thesis uses \emph{Renaissance}.
Renaissance is a semantic-level \Cpp\ analyser already in use at TNO-ESI.
It extracts the static call graph and populates the property graph
(\cref{sec:pg-decision}).
Renaissance resolves the full \ac{AST}, including template instantiations,
overload candidates, and macro expansions. It emits the property graph
schema described in \cref{sec:pg-schema}.

\subsection{Limitations in Event-Driven Systems}

Event-driven systems that route inter-component communication through
callback registration or publish-subscribe indirection pose a specific
challenge for static call-graph construction.
In CPSCore, subscriptions are registered via the IDC layer:
\begin{minted}{cpp}
IDC::subscribeOnPacket(channelId, handlerSlot);
\end{minted}
Standard \ac{AST} traversal identifies the call to
\texttt{subscribeOnPacket} but does not resolve which component will
receive a packet published on the same channel.
The consequence is a static call graph that contains edges \emph{into} the
IDC/Redis layer but lacks the cross-component routing information required
for interaction model reconstruction.
Resolving this limitation is part of the contribution of this thesis
(see \cref{sec:macro-resolution}).

\section{Execution Trace Analysis}\label{sec:trace-analysis}

An \emph{execution trace} is a temporally ordered sequence of events recorded
at runtime. Systems under analysis are typically instrumented with logging
or profiling statements to produce these events \cite{zaidman:traces}.
In this thesis, traces take the form of timestamped log entries that record
function entry and exit events. Each entry also records the component identity
and the topic name on the message bus.

\subsection{Challenges with Incomplete Traces}

Practical execution traces are almost always \emph{incomplete}:
\begin{itemize}
  \item Not all components are instrumented.
  \item Recording overhead may cause log entries to be dropped.
  \item A single scenario recording covers only a subset of the system's
    behaviour.
\end{itemize}

These gaps motivate the hybrid approach. The static call graph provides
structural completeness. The trace provides scenario specificity.

\section{Dependence-Based Slicing}\label{sec:slicing}

Program slicing, introduced by Weiser \cite{weiser:slicing}, computes the
subset of a program that affects the values of a set of \emph{criterion}
variables at a given point.
\emph{Dependence-based slicing} operates on the \ac{PDG}. The \ac{PDG}
combines the \ac{CDG} and the \ac{DDG} into a single graph \cite{ferrante:pdg}.

In the context of interaction model reconstruction, the slicing criterion is
the entry point of the scenario of interest (for example, a
\texttt{SynchronizedRunnerMaster::runStage} timeout event).
The slice then contains only those functions and message exchanges that are
reachable from the criterion node via dependence edges. This effectively filters
out interactions that belong to other scenarios running concurrently in the
system.

The adaptation used in this thesis (\cref{sec:fusion-slicing}) replaces the
traditional \ac{PDG} traversal with a scope-constrained graph query.
It operates over the property graph rather than over a source-level program
representation.

\section{Publish-Subscribe via Callback Registration}\label{sec:pubsub-background}

Many event-driven \Cpp\ systems decouple component communication through a
publish-subscribe pattern. In these systems, subscriptions are registered at
runtime via callback slots rather than direct function calls.
As a result, the publish-subscribe topology is \emph{implicit} in the
source code. It cannot be recovered by reading function call sites alone.

CPSCore implements this pattern through its \ac{IDC}
layer. This layer delegates to a pluggable \texttt{INetworkLayer} backed by
\texttt{boost::signals2} and Redis channels (\cref{sec:idc}).
Recovering the subscription topology requires matching publisher and subscriber
call sites by their shared channel identifier. This is the approach taken in
this thesis (see \cref{sec:macro-resolution}).
